\documentclass[12pt,a4paper]{article}
\usepackage[utf8]{inputenc}
\usepackage[vietnamese]{babel}
\usepackage[left=2cm,right=2cm,top=2cm,bottom=2cm]{geometry}
\usepackage{mathptmx} % Font Times New Roman
\usepackage{enumitem}

\pagestyle{plain}
\linespread{1.15}

\begin{document}

% --- HEADER QUỐC HIỆU & CƠ QUAN BAN HÀNH ---
\noindent
\begin{minipage}[t]{0.45\textwidth}
    \begin{center}
        \small
        CÔNG AN THÀNH PHỐ HÀ NỘI\\
        \textbf{\underline{PHÒNG CẢNH SÁT HÌNH SỰ}}\\[0.1cm]
        Số: 07c/BB-LK\\
        \textit{V/v lấy lời khai người bị nghi vấn}
    \end{center}
\end{minipage}
\hfill
\begin{minipage}[t]{0.52\textwidth}
    \begin{center}
        \small
        \textbf{CỘNG HÒA XÃ HỘI CHỦ NGHĨA VIỆT NAM}\\
        \textbf{\underline{Độc lập - Tự do - Hạnh phúc}}\\[0.1cm]
        \textit{Hà Nội, ngày 25 tháng 07 năm 2026}
    \end{center}
\end{minipage}

\vspace{0.6cm}

% --- TIÊU ĐỀ VĂN BẢN ---
\begin{center}
    \textbf{\large BIÊN BẢN LẤY LỜI KHAI}\\
    \textit{(Vụ án mạng tại số 14 Đường Bờ Sông, Phân khu Cảng)}
\end{center}

\vspace{0.3cm}

\noindent Vào hồi 18 giờ 00 phút, ngày 25 tháng 07 năm 2026, tại Phòng Cảnh sát Hình sự Công an TP. Hà Nội.\\
Chúng tôi gồm:
\begin{enumerate}[label=\arabic*., leftmargin=1.5cm]
    \item Điều tra viên: Đại úy Lê Minh — Trưởng ban chuyên án.
    \item Cán bộ ghi biên bản: Trung úy Nguyễn Văn Hoàng — Cán bộ Đội Án mạng.
\end{enumerate}

\noindent Tiến hành lấy lời khai người bị nghi vấn:
\begin{itemize}[leftmargin=1.5cm]
    \item \textbf{Họ và tên:} NGUYỄN THANH TÙNG. Giới tính: Nam. Sinh ngày: 18/03/1991.
    \item \textbf{CCCD số:} 001091001823 cấp ngày 12/01/2022 tại Cục CSQLHC về TTXH.
    \item \textbf{Nơi ĐKHKTT:} KTT Nhà máy Sứ, Phân khu Cảng, TP. Hà Nội.
    \item \textbf{Nghề nghiệp:} Tự do. Mối quan hệ với nạn nhân: Bạn từ thời thơ ấu.
\end{itemize}

\section*{NỘI DUNG LỜI KHAI (HỎI VÀ ĐÁP)}

\noindent \textbf{Hỏi (ĐTV Lê Minh):} Anh có mặt tại nhà nạn nhân Khang lúc mấy giờ tối 24/07/2026 và đã xảy ra chuyện gì?\\
\textbf{Đáp (Nguyễn Thanh Tùng):} Tôi... tôi xin đầu thú sự việc! Tối qua khoảng 20:00 tôi đến gặp Khang để xin lại mảnh giấy note kỷ vật trò trốn tìm năm 1998 (ngày em trai tôi bị ngạt thở trong tủ). Khang không cho mà còn buông lời nhạo báng cái chết của em tôi và đe dọa tống tiền ép tôi gánh nợ. Tôi quá phẫn uất lao vào giằng co, đập vỡ bộ bình cốc thủy tinh trên bàn trà.

\noindent \textbf{Hỏi (ĐTV Lê Minh):} Anh đã ra tay gây thương tích cho nạn nhân như thế nào?\\
\textbf{Đáp (Nguyễn Thanh Tùng):} Trong lúc giằng co, tôi dùng hết sức xô mạnh Khang ra. Khang bị ngã ngửa, đầu đập mạnh vào mép gỗ tủ âm tường phát ra tiếng "bịch" rồi ngất xỉu gục dưới sàn nhà. Thấy Khang gục ngất im lìm và máu rỉ ra từ gáy, tôi quá hoảng sợ tưởng Khang đã chết nên tháo chạy vội vã ra lấy xe máy phóng đi lúc hơn 20:15!

\noindent \textbf{Hỏi (ĐTV Lê Minh):} Anh có dùng mảnh thủy tinh vỡ đâm vào cổ nạn nhân không?\\
\textbf{Đáp (Nguyễn Thanh Tùng):} Không! Tôi thề với trời đất! Tôi chỉ xô Khang ngã gục ngất xỉu rồi hoảng quá bỏ chạy ngay lúc 20:15. Tôi tuyệt đối không dùng mảnh thủy tinh nào đâm Khang cả!

\noindent Biên bản lấy lời khai kết thúc vào hồi 19 giờ 30 phút cùng ngày. Biên bản này đã được đọc lại cho anh Nguyễn Thanh Tùng nghe, công nhận đúng và cùng ký tên xác nhận.

\vspace{0.8cm}

% --- CHỮ KÝ NGHỊ ĐỊNH 30 ---
\noindent
\begin{minipage}[t]{0.45\textwidth}
    \small
    \textbf{\textit{Nơi nhận:}}\\
    - Hồ sơ vụ án TEST-99;\\
    - Viện KSMND TP. Hà Nội;\\
    - Lưu: VT, CSHS.
\end{minipage}
\hfill
\begin{minipage}[t]{0.5\textwidth}
    \begin{center}
        \small
        \textbf{NGƯỜI KHAI BÁO}\\
        \textit{(Ký, ghi rõ họ tên)}\\[1.5cm]
        \textbf{Nguyễn Thanh Tùng}
    \end{center}
\end{minipage}

\end{document}
